% /root/monodromy, 19 November 2026. Complete root and reduction/descent audits passed.
\subsection{Uniform degree-one image for every fixed mixed product}
\label{subsec:mixed-uniform-degree-one-inertia-candidate}
\status{Proved and independently checked. The root and reduction/descent
agents audited all six complete assertions, the source dependencies and
the full-image assembly. Both independently read and reran the universal
Pfaffian and stable-loss certificates; the rank-compatibility certificate
was independently read and rerun by the author and reduction/descent agent.
The result covers every ordinary modulus.}

Work over an algebraically closed field $k$ of characteristic two. For
$a\in k^\times$ put
\[
 E_a:y^2+xy=x^3+a,\qquad E_0:z^2+z=w^3,
 \qquad X_a=E_a\times E_0,
 \qquad L=\mathcal O_{E_a}(3O)\boxtimes\mathcal O_{E_0}(3O).
\]
Write $W=H^0(X_a,L)$, $P=\mathbf P(W)$ and $s=(s_{ij})$ in
the tensor basis $(1,x,y)\otimes(1,w,z)$; put $D(s)=\det(s_{ij})$.
The chosen invariant volume is $(dx/x)\wedge dw$.

\begin{proposition}[Two cubic components at every ordinary modulus]
\label{prop:mixed-uniform-degree-one-Hasse-candidate}
Let
\[
 R_o=\begin{pmatrix}1&0&a\\0&0&a\\0&a&a\end{pmatrix},
 \qquad J=(\delta_{i,2-j})_{0\le i,j\le2},
 \qquad R=R_o\otimes J.
\]
The relative coherent Frobenius bundle is $\Omega_P^1$, and its
determinant is, up to a nonzero constant,
\[
                   h(s)=D(s)^2Q(s),\qquad Q(s)=s^tRs^{[2]}.
\]
The cubics $D,Q$ are geometrically irreducible and distinct, and
$Q$ is smooth. The generic relative stable rank is eight. At the
geometric generic point of $Q$, it is seven, and the linearized
relative Frobenius has a one-dimensional kernel.
\end{proposition}

\begin{proof}
The root-removed Cartier multiplication matrices on the two factors are
\[
 M_o(u)=\begin{pmatrix}
 u_0&a u_2&a u_1\\u_2&u_1&u_0\\0&u_2&u_1+u_2
 \end{pmatrix},\qquad
 M_0(u)=\begin{pmatrix}u_1&u_0&0\\u_2&0&u_0\\0&u_2&u_1\end{pmatrix}.
\]
For the first factor, the underlying Cartier identities are
$C(1)=1$, $C(x)=0$, $C(y)=x$, and $C(xy)=y+\sqrt a$,
together with $C(f^2)=f$ for an invariant Cartier-fixed
differential. The second-factor identities were proved in
Proposition~\ref{prop:ss-product-relative-Hasse}.
If $M_{o,i},M_{0,j}$ are the coefficient matrices, put
$M(s)=\sum s_{ij}M_{o,i}\otimes M_{0,j}$.
The symmetric multiplication identity gives $M(s)s=0$.
More explicitly, $A(s)=RM(s)$ is alternating, and its coefficient
tensor is symmetric in all three arguments. One may check the latter
directly from
\[
 R_oM_o(u)=\begin{pmatrix}
 u_0&0&a u_2\\
 0&a u_2&a(u_1+u_2)\\
 a u_2&a(u_1+u_2)&a(u_0+u_1+u_2)
 \end{pmatrix}
\]
and the alternating volume tensor on the second factor.
The nine size-eight Pfaffians satisfy
\[
                    \operatorname{Pf}_8(A(s))=a^3D(s)s,
                    \qquad \det R=a^6.
\]
These are polynomial identities over $\mathbf F_2[a,s_0,\ldots,s_8]$:
the retained exact certificate expands every Pfaffian recursively and
checks all nine identities. The alternating adjugate identity consequently
gives
\[
                  \operatorname{adj}M(s)=D(s)^2s\,s^tR.
\]
The universal hypersurface sequence is the Euler sequence, since
$H^1(X_a,L^{-1})=0$ and $h^2(X_a,\mathcal O_{X_a})=1$.
Frobenius on the last one-dimensional quotient is zero, because the
surface is nonordinary. Thus the relative map is the restriction of
$M(s)^t$ from $(s^{[2]})^\perp$ to $s^\perp$. Its cofactor formula
is exactly $D(s)^2s^tRs^{[2]}$, as in the preceding supersingular
calculation. Its degree is nine, agreeing with
$\det\Omega_P^1=\mathcal O_P(-9)$.

The determinant cubic $D$ is irreducible by its matrix-factorization
parametrization. Since
$\nabla Q=Rs^{[2]}$ and $R$ is invertible, $Q$ is smooth and hence
irreducible. The coefficient vector with $s_{01}=1$ and all other
entries zero has $D=0,Q=1$, so the two factors are distinct. Their
complement is a nonempty full-rank open, proving generic rank eight.

We verify stable loss one on $Q$ uniformly in $a$, not merely at a
single ordinary modulus. Let $z$ be an auxiliary nonzero parameter
and take, in lexicographic order,
\[
             s=(z^3/a,z,0,0,1,0,0,0,1).
\]
Then $D=z^3/a\ne0$ and $Q=0$. The vector
\[
 u=\left(0,\frac{z^3}{a},\frac a z,0,
          \frac{a+z^3}{z^3},\frac{a^2+z^5}{az^2},0,0,\frac a{z^3}\right)^t
\]
satisfies $M(s)u=s^{[2]}$, and
\[
             s^tR u^{[2]}=\frac{a^4+a^3z^4+z^{10}}{a^2z^3}.
\]
Both identities are also checked over $\mathbf F_2(a,z)$ by the
retained program. The last expression is not identically zero as
a function of $z$ for any $a\ne0$.
On $Q\setminus D$ the full matrix has kernel $ks$, and the
relative Cartier kernel is the line generated by $u$ modulo $ks$.
Its squared representative does not lie in $\operatorname{im}M(s)$
precisely when the displayed pairing is nonzero, since $s^tR$
spans the left kernel and $s^{[2]}$ already lies in the image.
Thus the first relative image has rank seven and is already
Frobenius-stable at such a point. Openness of the corresponding
iterate minor proves the generic assertion on the irreducible $Q$.
\end{proof}

\begin{lemma}[Full elementary inertia from a moving kernel line]
\label{lem:moving-kernel-full-root-inertia-candidate}
Let $Y$ be a smooth variety over a perfect field of characteristic two,
and let $D\subset Y$ be an integral divisor. Let a vector bundle
$\mathcal E$ of rank $r$ carry regular semilinear Frobenius. Suppose
its geometric generic rank is $r$, its stable rank at the generic point
of $D$ is $r-1$, its linearized kernel there is a line, and its
determinant has order one along $D$.
Suppose that on a neighbourhood of the generic point of $D$ there is
an inclusion $\mathcal E\subset V\otimes\mathcal O_Y$ as a subbundle of a constant
bundle, and the linearized kernel line in $F^*\mathcal E$ defines
a rational map
\[
                      \kappa:D\dashrightarrow\mathbf P(V^{(2)})
\]
whose differential has rank $r-1$ at the generic point.
Then the divisorial inertia image of the generic Lang local system
contains the entire elementary group
\[
           \operatorname{Hom}(\mathbf F_2,\mathbf F_2^{r-1}),
\]
acting trivially on the regular $(r-1)$-dimensional unit subspace
and on the one-dimensional quotient. In particular it contains
every transvection with that fixed hyperplane.
\end{lemma}

\begin{proof}
Pass to the strict henselization at the generic point of $D$.
Write $R$ for this discrete valuation ring, $t$ for a uniformizer,
and $K$ for its separably closed residue field. The finite part of
the Lang scheme supplies $r-1$ regular Frobenius-fixed sections
$w_1,\ldots,w_{r-1}$, with independent special values spanning
the stable part. This is the finite-part construction already
proved using Stacks~\cite{Stacks0BTN,Stacks09ZS}. Complete them
to a frame by $e$. In that frame
\[
                   F=\begin{pmatrix}I_{r-1}&b\\0&a\end{pmatrix},
                  \qquad a=t c,\qquad c\in R^\times.
\]
No perfectness of $K$ is assumed. In particular the constant
column $b_0=b\bmod t$ need not be removable by a regular change
of frame. The linearized kernel is the line generated by
\[
             e^{[2]}+\sum_{i=1}^{r-1}b_{0,i}w_i^{[2]}.
\]
Use the given inclusion in the constant bundle. In Kähler
differentials over the perfect ground field, every entry of
$e^{[2]}$ and $w_i^{[2]}$ has zero differential. The differential
rank of the projective kernel line is therefore the rank of
$db_{0,1},\ldots,db_{0,r-1}$. The assumed rank $r-1$ makes
these differentials linearly independent over $K$. This calculation
is unchanged by the separable residue-field extension: separable
algebraic extension preserves the differential rank.

The nonzero fixed vector of the quotient has coordinate $a^{-1}$.
Lifting it requires the independent Artin--Schreier equations
\[
                      y_i^2+y_i=b_i/a^2.
\]
For any nonzero $(\epsilon_i)\in\mathbf F_2^{r-1}$, the leading
coefficient of $\sum\epsilon_i b_i/a^2$ at order $t^{-2}$ is
$c_0^{-2}\sum\epsilon_i b_{0,i}$. Its differential is nonzero,
so it is not a square in $K$. If this sum were $y^2+y$ in the
fraction field of $R$, its valuation would force $v_t(y)=-1$,
whose leading coefficient would make that coefficient a square.
This is impossible. The same argument applies in the henselian
field itself and does not require completion.

The $r-1$ classes are consequently linearly independent in
$\operatorname{Frac}(R)/\{y^2+y\}$. Successive quadratic
Artin--Schreier extensions therefore have compositum Galois group
$\mathbf F_2^{r-1}$. This standard elementary assertion follows
also by taking the annihilator of the image in
$\mathbf F_2^{r-1}$: a proper image would give exactly one of the
excluded nonzero linear combinations. Over a strictly henselian
trait this is its inertia image, including possible purely
inseparable residue extensions. Its action adds the arbitrary
constant combination $\sum\epsilon_iw_i$ to one chosen lift of
the quotient vector, and fixes all the $w_i$. This is exactly
the asserted full elementary group.
\end{proof}

\begin{corollary}[Seven independent root directions for every mixed modulus]
\label{cor:mixed-every-modulus-full-root-inertia-candidate}
For every $a\ne0$, the geometric degree-one relative image for
$X_a,L$ contains the group of order $2^7$ consisting of all
transvections with a fixed seven-dimensional hyperplane, together
with the identity.
\end{corollary}
\begin{proof}
Apply the preceding lemma to the cubic $Q$ on $P=\mathbf P^8$.
The stable rank and simple determinant order were just proved.
On an affine coefficient chart, the relative Euler bundle embeds
in the constant nine-dimensional bundle. The kernel of the
linearization on $Q\setminus D$ is the line
$kR^ts\subset (s^{[2]})^\perp$. Thus its projective map is
the restriction of the projective linear isomorphism
$[s]\mapsto[R^ts]$. Its differential on the smooth cubic has
rank seven. All hypotheses hold. The inertia is geometric,
since the construction was over an algebraically closed field.
Restricting the coefficient invertibility open to its dense
smooth-curve open preserves the finite image, as before.
\end{proof}

\begin{lemma}[The exceptional ranks of the mixed Cartier tensor]
\label{lem:mixed-degree-one-exceptional-ranks-candidate}
For the alternating matrix $A(w)$ above put
\[
       \mathcal Z=\{(s,v):v\ne0,\ A(R^{-t}v^{[2]})s=v\}.
\]
The part of $\mathcal Z$ where the $3\times3$ matrix $v$ has
rank at most two has dimension at most eight.
\end{lemma}
\begin{proof}
Write $B(u)=R_oM_o(u)$. Directly,
\[
 \det B(u)=a^2\bigl(u_0^2u_2+u_0u_1u_2+u_0u_1^2+a u_2^3\bigr).
\]
The cubic in parentheses is smooth for every $a\ne0$.
Its three partial derivatives are
$u_1^2+u_1u_2$, $u_0u_2$, and
$u_0^2+u_0u_1+a u_2^2$.
If $u_2=0$, their simultaneous vanishing forces $u_1=u_0=0$;
if $u_2\ne0$, the second forces $u_0=0$ and the third is
nonzero. Thus the projective cubic is smooth and contains no
line. In particular $B(u)$ has rank three off that cubic
and rank two on it: rank at most one would make all cofactors,
and therefore all three partial derivatives, zero.

Suppose first that $v$ has matrix rank two. Write
\[
 v=t_1\otimes q_1+t_2\otimes q_2,
 \qquad T=(t_1,t_2),\qquad Q_0=(q_1,q_2),
\]
where both displayed $3\times2$ frames have full rank.
All such representations form a $\GL_2$-torsor of dimension
four over the rank-two matrix locus. Put
\[
 u_i=R_o^{-t}t_i^{[2]},\qquad r_i=Jq_i^{[2]},\qquad
 w=R^{-t}v^{[2]}=u_1\otimes r_1+u_2\otimes r_2.
\]
The plane spanned by $u_1,u_2$ cannot lie in the smooth cubic
$\det B=0$. A change of the two-element frame therefore makes
$B(u_1)$ invertible; square roots needed for that change exist
over the algebraically closed field. Thus the open of such
frames meets every rank-two orbit.

Contract the second tensor factor with $r_1,r_2$, and complete
them to a basis. The alternating volume tensor on that factor
shows that a vector lies in $\operatorname{im}A(w)$ precisely
when its first two contractions can be written
$B(u_2)z,B(u_1)z$ for one vector $z$. The remaining contraction
is unrestricted, because $B(u_1)$ is invertible. In particular
$A(w)$ has rank six. Define
\[
 H=(h_{ij})_{1\le i,j\le2},\qquad
 h_{ij}=q_i^tJq_j^{[2]},\qquad C=B(u_2)B(u_1)^{-1}.
\]
The required compatibility is exactly
\[
                         T H_{\bullet1}=CT H_{\bullet2}.
\]
For fixed $T$, these are $k_T=\dim\langle T,CT\rangle$
independent linear equations on the four entries of $H$,
where $k_T$ is two or three.

The Gram map $Q_0\mapsto Q_0^tJQ_0^{[2]}$ on full-rank
$3\times2$ frames is smooth of relative dimension two.
Indeed its derivative is
\[
                     dH=(dQ_0)^tJQ_0^{[2]},
\]
which is surjective onto all $2\times2$ matrices: each row
of $(dQ_0)^t$ may be chosen independently and
$JQ_0^{[2]}$ has rank two. Hence the compatible $Q_0$ have
dimension at most $6-k_T$.

The exceptional condition $k_T=2$ is a proper closed condition
on the six-dimensional open of $T$-frames under consideration.
To check properness uniformly, take $T=(e_0,e_1)$. Then
\[
 u_1=(1,1,0)^t,\quad u_2=(0,a^{-1},a^{-1})^t,\qquad
 B(u_1)=\begin{pmatrix}1&0&0\\0&0&a\\0&a&0\end{pmatrix},
 \quad B(u_2)=J,
\]
and $Ce_0=e_2\notin\langle e_0,e_1\rangle$.
The locus with $k_T=3$ therefore has compatible-frame
dimension at most $6+(6-3)=9$; the locus with $k_T=2$ has
dimension at most $5+(6-2)=9$. Subtracting the four-dimensional
frame fibre gives a compatible $v$-locus of dimension at most
five. The affine space of solutions $s$ has dimension
$9-\rank A(w)=3$, giving the required bound eight.

For rank one write $v=t\otimes q$, with $t,q\ne0$, and put
$u=R_o^{-t}t^{[2]}$, $r=Jq^{[2]}$.
Then $A(w)$ is the tensor product of $B(u)$ and the
rank-two alternating volume matrix determined by $r$.
If $B(u)$ is invertible, the rank is six; the entire rank-one
$v$-locus already has dimension five, so its incidence has
dimension at most $5+(9-6)=8$.
If $B(u)$ has rank two, the rank of $A(w)$ is four.
Compatibility necessarily requires
\[
                             q^tJq^{[2]}=0,
\]
since the image on the second factor is the plane annihilating
$r$. The condition on $u$ cuts the nonzero $t$-space to
dimension two, and the displayed nonzero cubic cuts the
nonzero $q$-space to dimension two. Quotienting their
one-dimensional scaling ambiguity gives dimension at most
three for compatible $v$. Adding the solution dimension five
again gives eight. These cases exhaust all nonzero $v$ of
matrix rank at most two.
\end{proof}

\begin{lemma}[Transitivity on all nonzero relative vectors]
\label{lem:mixed-degree-one-Lang-transitive-candidate}
For every $a\ne0$, the geometric degree-one relative image is
transitive on the $255$ nonzero vectors of its eight-dimensional
natural representation. In particular it is irreducible.
\end{lemma}
\begin{proof}
Work on the affine coefficient cone where $D(s)Q(s)\ne0$.
The nonzero Lang scheme is finite etale of degree $255$ there.
In the ambient nine-dimensional frame its equations are
\[
                    v=M(s)^tv^{[2]},\qquad s^tv=0,\qquad v\ne0.
\]
Put $w=R^{-t}v^{[2]}$. Full symmetry of the coefficient tensor
gives $A(s)w=A(w)s$, so the first equation is equivalent to
$A(w)s=v$. The second equation then holds automatically,
because $s^tA(w)s=0$.

When $v$ has matrix rank three, $D(w)\ne0$ and
$\ker A(w)=kw$. The equation is solvable precisely when
\[
                         v^tR^{-t}v^{[2]}=0.
\]
This is a smooth projective cubic, since its gradient is
$R^{-t}v^{[2]}$; its punctured affine cone is irreducible.
Its full-matrix-rank open is nonempty: its monomials have a
squared variable, while every determinant monomial has three
distinct variables, so it is not the determinant cubic.
Over that open the solution space in $s$ is an affine-line
bundle, hence irreducible of dimension nine.

The previous lemma bounds the remaining incidence by dimension
eight. Every component of a finite etale cover of the integral
nine-dimensional coefficient cone has dimension nine and
surjects onto that cone. Thus none can lie in the exceptional
incidence. All meet the one irreducible affine-line incidence,
so the nonzero Lang cover is irreducible. Its degree is $255$,
which proves the asserted transitivity. Pullback to the affine
cone can only decrease the monodromy image, so the same holds
on the projective parameter open and then on its dense
smooth-curve open. A nonzero proper invariant subspace would
have a proper nonempty invariant set of nonzero vectors,
contradicting transitivity.
\end{proof}

\begin{theorem}[Full degree-one image for every mixed product]
\label{thm:mixed-every-fixed-modulus-degree-one-candidate}
Let $k$ be any perfect field of characteristic two, let $E/k$
be an ordinary elliptic curve and let $E'/k$ be a supersingular
elliptic curve. On
\[
 X=E\times E',\qquad
 L=\mathcal O_E(3O)\boxtimes\mathcal O_{E'}(3O),
\]
the relative Hasse--Witt system on the maximal-rank smooth
open of $|L|$ has rank eight and
\[
                         G_{\mathrm{geom}}=G_{\mathrm{arith}}
                              =\GL_8(\mathbf F_2).
\]
This holds for every ordinary modulus, including every algebraic
one. Its geometric representation is irreducible.
\end{theorem}
\begin{proof}
Over an algebraic closure the pair has the displayed form
$E_a\times E_0$ for some $a\ne0$. Let $G$ be its geometric
image on $V=\mathbf F_2^8$. The preceding lemma proves
irreducibility. The inertia corollary supplies, for a nonzero
$\phi\in V^*$, the entire group
\[
                  U_\phi=\{1+u\otimes\phi:\phi(u)=0\}.
\]
The span of the $G$-orbit of $\phi$ is all of $V^*$, by dual
irreducibility. Choose a basis $\phi_1,\ldots,\phi_8$ from
that orbit and let $e_1,\ldots,e_8$ be its dual basis.
Conjugation gives $U_{\phi_i}\subset G$ for every $i$.
Consequently $1+e_j\otimes\phi_i\in G$ whenever $i\ne j$.
These are all elementary matrices in that basis and generate
$\SL_8(\mathbf F_2)=\GL_8(\mathbf F_2)$. This closure uses
only elementary matrices, without a finite-group classification.
The arithmetic group contains the geometric group and is
contained in the same general linear group, proving equality.
Geometric base extension preserves connectedness of the finite
frame cover, so the conclusion holds over every stated perfect
field and for all twists of the elliptic factors.
\end{proof}

\begin{remark}[Retained computations and higher-degree scope]
\label{rem:mixed-uniform-inertia-scope-candidate}
The exact programs
\path{agents/monodromy/code/verify_mixed_degree_one_pfaffians.py}
and
\path{agents/monodromy/code/explore_mixed_universal_Q_pencil.py}
retain all universal polynomial and rational identities used above.
Earlier failed point choices are preserved separately and are not
inputs to this proof. The parent independently supplied and checked
the exceptional-rank dimension argument. Its exact auxiliary certificate is
\path{code/mixed_lang_rank_compatibility.py}, with retained output
\path{outputs/mixed_lang_rank_compatibility.json}; it checks the
cubic determinant and partials, the uniform exceptional-frame witness,
the Gram differential and its block-minor identities. These finite
identities support the full dimension proof above.
The theorem closes the fixed-modulus degree-one case. Higher degrees
at an arbitrary fixed algebraic ordinary modulus require a further
uniform argument; the already proved generic-modulus all-degree
theorem remains a separate statement.
\end{remark}
